\subsection{What the demonstration distribution can teach}

SFT can teach instruction following, answer format, tool syntax, terminology, length, and
reasoning style when these behaviors are represented consistently in the demonstrations.
Its generalization is constrained by what varies in that dataset.

\paragraph{System prompts as conditions.}
Including a system prompt in the input does not guarantee that the learned behavior
remains controllable through that prompt. If exactly the same system message appears in
every example, it is a constant rather than an informative variable. The data provide no
contrast showing which behavior should change when the message changes, so repeated SFT
can absorb its instruction into the model parameters. There is no universal iteration at
which this happens, and a larger training run does not solve the identifiability problem.
Prompt control is better tested and taught by varying system instructions, including
contrasting instructions for comparable user inputs, sometimes omitting generic prompts,
and evaluating behavior when the prompt is changed or removed.

\paragraph{One demonstrated path.}
A demonstration selects one valid token path. Cross-entropy increases that path's local
probability even when other continuations would also be correct. Repetitive data can
therefore reduce diversity, reinforce template artifacts, and shift capability away from
domains absent from the mixture. Low training loss establishes imitation of the training
set, not generalization.

\paragraph{Teacher-forced prefixes.}
During SFT, the model predicts each response token after the correct demonstrated prefix.
During inference, described in Section~\ref{sec:transformer-inference}, each new token is
conditioned on the model's own earlier samples. An early mistake can lead to a prefix that
never appeared in the SFT data. This exposure-bias gap explains why dense demonstration
learning can produce strong initial behavior without directly teaching recovery from the
model's own errors.

SFT is therefore a strong initialization method when desired behavior can be demonstrated
clearly. It does not by itself identify every valid alternative, evaluate terminal
outcomes, or supply credit for actions taken on the model's own trajectory. Those are
separate problems for the later sections on reinforcement learning, critics, and
on-policy distillation.
